\section{Verification Time of the Alipay Models}
\label{app:alipay-time}

The three Alipay models of Xiao et al.~\cite{xiao2026payment} contain a low-entropy payment password 
but no function output from which that password could be recomputed, 
so they carry a low-entropy value without carrying a reachable brute-force attack.
They therefore measure what the translation costs when there is nothing for the oracles to find.
Table~\ref{tab:perf-baseline} places each annotated model beside the same model without annotations.
Every lemma reaches the same verdict in both versions and the finished proofs have the same number of steps, 
so the annotation changes the cost of the analysis and not its outcome.

For the application and mobile-web models the annotation adds about twelve seconds, 
under nine percent of the running time.
The order-code model is more expensive in both versions.
Its original tactic leaves attacker-knowledge goals about session values ahead of everything else,
and adding a single deprioritisation line brings the unannotated model to 42\% and the annotated
model to 32\% of its unadjusted time.
We report both variants, and add that line to both versions of the model, so that the pair being
compared differs only in the annotation.
The remaining gap of about four minutes is what the translation costs on this model.

\begin{table}[t]
    \centering
    \caption{Verification time of the Alipay models with and without low-entropy annotations, under \Tamarin{}~1.10.0 and Maude~3.1. ``Adjusted'' adds one deprioritisation line to the tactic of both versions.}
    \label{tab:perf-baseline}
    \small
    \renewcommand{\arraystretch}{1.15}
    \begin{tabularx}{\columnwidth}{@{}>{\raggedright\arraybackslash}X r r r@{}}
    \toprule
    \textbf{Model} & \textbf{Plain (s)} & \textbf{EntropyVerif (s)} & \textbf{$\Delta$ (s)} \\
    \midrule
    Alipay, application          & 134.41 & 146.03 & $+11.62$ \\
    Alipay, mobile web           & 134.10 & 146.16 & $+12.06$ \\
    Alipay, order code           & 456.57 & 1384.68 & $+928.11$ \\
    Alipay, order code, adjusted & 191.08 & 438.45 & $+247.37$ \\
    \bottomrule
    \end{tabularx}
\end{table}

\section{Where the Proof Effort Goes}
\label{app:breakdown}

This appendix gives a detailed breakdown of the proof effort, using the BC case study as the example, 
and reports how much of a completed proof is spent on the goals our translation
introduces.
\Tamarin{} solves constraints in an interdependent way, so its running time cannot be divided
among individual goals, and we decompose the finished proofs instead.
We use BC because it is the only case study with verified all-traces lemmas, 
which are the proofs where the search has to be exhausted and the cost of the translation therefore shows.

\subsection{Method}

\Tamarin{} reports a step count for each lemma.
In the proof script that count corresponds to the \Tterm{case} lines: a \Tterm{solve} on a goal
splits into one or more cases, 
siblings are separated by \Tterm{next}, and the goal is discharged by \Tterm{qed}.
For the \Tterm{SecrecyCentralSK} lemma of one configuration, for instance, the script contains
3292 \Tterm{solve} lines, 3292 \Tterm{qed} lines, 6359 \Tterm{next} lines 
and 9651 \Tterm{case} lines, and \Tamarin{} reports 9653 steps.
Counting cases therefore reproduces the reported figure, and we take one case to be one step.
We assign each step to the goal of the innermost enclosing \Tterm{solve}, and read the kind of that goal from its fact:
\TFact{MDerive0}, \TFact{MDerive1} and \TFact{EntropyLow} are the facts our translation adds, 
\TFact{KU} is attacker knowledge, and the remaining facts belong to the protocol model.
We also record which rule closes each step.

We restrict the count to the lemmas \Tamarin{} verified.
A falsified lemma stops at the first counterexample it finds, so its proof is a single path and
its step count measures how quickly an attack was reached rather than the cost of a search.
A verified all-traces lemma is the opposite case: every branch has to be closed for the property to hold, 
so every case in the script is work that was actually carried out.
The \Tterm{Executability} lemma is also excluded, being an existence property that closes in thirteen to fifteen steps everywhere.
Eleven verified lemmas remain, 1{,}779{,}376 steps in total.

\subsection{Results}

Table~\ref{tab:perf-breakdown-kind} gives the split by kind of goal.
The facts our translation adds account for 86.4\% of the steps, while the goals of the protocol
model itself account for 54 steps out of 1{,}779{,}376.
The tactic sorts the protocol goals ahead of everything else, so they are settled early and the
remainder of the search runs on derivation facts and attacker knowledge.

\begin{table}[t]
    \centering
    \caption{Proof steps of the eleven verified BC lemmas, by kind of goal.}
    \label{tab:perf-breakdown-kind}
    \small
    \renewcommand{\arraystretch}{1.15}
    \begin{tabularx}{\columnwidth}{@{}>{\raggedright\arraybackslash}X r r@{}}
    \toprule
    \textbf{Goal} & \textbf{Steps} & \textbf{Share} \\
    \midrule
    \TFact{MDerive0}          & 1495256 & 84.0\% \\
    \TFact{MDerive1}          &   42504 &  2.4\% \\
    \TFact{KU}                &  241562 & 13.6\% \\
    Protocol state and action &      54 &  0.0\% \\
    \midrule
    Total                     & 1779376 & \\
    \bottomrule
    \end{tabularx}
\end{table}

The \TFact{MDerive0} steps are the largest group, and the rules that close them show what they
are: 72\% are closed by a rule of the BC model, led by \Tterm{PerLegacyEncRsp} with 265{,}936
steps, \Tterm{PerLegacyRecvAuRand} with 153{,}552 and \Tterm{PerSendKeysizeReq} with 139{,}216,
while the remaining 28\% are closed by the two derivation rules the compiler emits.
A \TFact{MDerive0} goal asks which terms an observed message was computed from, and \Tamarin{}
answers it by trying the protocol rules that could have produced that message.
The cost of the translation therefore grows with the number of protocol rules that can produce an
observed message, not with the number of keys marked as low entropy.

The \TFact{KU} steps are spent rebuilding key material.
Constructing the BC key-derivation functions accounts for most of them, \Tterm{c\_h4} alone for
118{,}826 steps, or 49\% of the group, followed by \Tterm{c\_combine} with 69{,}316 steps, or
29\%, which assembles a key from separately recovered parts.
Propagation is minor: the \TFact{MDerive1} facts, which the propagation rule produces, take 42{,}504
steps, one thirty-fifth of the \TFact{MDerive0} total.

The rule that performs the recovery itself, \Tterm{bf\_low\_entropy\_message}, closes 4{,}980
steps, which is 0.28\% of the proofs.

Table~\ref{tab:perf-breakdown} gives the same split per lemma.
Within one configuration the proportions are close to constant across lemmas, so what the
translation costs is a property of the configuration rather than of a particular security
property.

\begin{table}[t]
    \centering
    \caption{Proof steps per verified lemma. MD0, MD1: the derivation facts our translation adds. KU: attacker knowledge. Prot.: goals of the protocol model.}
    \label{tab:perf-breakdown}
    \small
    \renewcommand{\arraystretch}{1.15}
    \begin{tabularx}{\columnwidth}{@{}>{\raggedright\arraybackslash}X l r r r r@{}}
    \toprule
    \textbf{Configuration} & \textbf{Lemma} & \textbf{MD0} & \textbf{MD1} & \textbf{KU} & \textbf{Prot.} \\
    \midrule
    \multirow{3}{*}{\makecell[l]{LgOn-Legacy $\times$\\LgOn-Enhanc, QA}}
        & SecCenSK & 169584 & 5010 &  2254 &  1 \\
        & SecPerSK & 261624 & 7998 &  4682 &  7 \\
        & MitM     & 169584 & 5010 &  2254 &  5 \\
    \midrule
    \multirow{3}{*}{\makecell[l]{ScLg-Secure $\times$\\ScLg-Enhanc, QA}}
        & SecCenSK & 133264 & 3000 & 30028 &  2 \\
        & SecPerSK & 221184 & 4980 & 52477 &  8 \\
        & MitM     & 266528 & 6000 & 60084 & 10 \\
    \midrule
    \multirow{3}{*}{\makecell[l]{ScOn-Secure $\times$\\ScOn-Enhanc, QA}}
        & SecCenSK &  10176 &  324 & 12226 &  1 \\
        & SecPerSK &  15632 &  486 & 19549 &  3 \\
        & MitM     &  10176 &  324 & 12240 &  5 \\
    \midrule
    \multirow{2}{*}{\makecell[l]{ScOn-Secure $\times$\\ScOn-Legacy, QA}}
        & SecCenSK &   7776 &  324 &  1550 &  1 \\
        & MitM     &   7776 &  324 &  1564 &  5 \\
    \midrule
    \multirow{2}{*}{\makecell[l]{ScOn-Secure $\times$\\ScOn-Legacy, SA}}
        & SecCenSK & 110976 & 4362 & 21320 &  1 \\
        & MitM     & 110976 & 4362 & 21334 &  5 \\
    \bottomrule
    \end{tabularx}
\end{table}

\subsection{Summary}

Three things follow.
The goals our translation adds take 86.4\% of the steps of a verified BC proof, and the protocol
model's own goals take 54 of 1{,}779{,}376.
Those steps are not new cryptographic reasoning but a second traversal of the protocol rules,
asking of each observed message which rule could have produced it, which is why 72\% of them are
closed by a protocol rule.
Recovering a key is almost free at 0.28\% of the steps and propagation is cheap at 2.4\%, so what
the analysis pays for is reaching the point at which a recovery becomes possible, and it pays most
where no attack exists and every branch has to be closed.
